\section{Conclusion}
\label{sec:conclusion}

We studied the preparation of $\plusL$ in the $XYZ^2$ hexagonal code by measurement, decoding and recovery under circuit-level noise. Three choices make the protocol fault tolerant. The data qubits start in a product state chosen in the CSS frame of the code, which fixes $\Xb$ and the $d^2$ independent generators of $\Sz$. All generators are measured with a depth-six schedule selected by an exhaustive search, and the logical operator is never measured directly. The full space-time syndrome is decoded once, and the recovery is a Pauli frame update. The resulting protocol has fault distance $d+1$, which is the largest possible for $\plusL$. With one round of measurements the threshold of matching under depolarizing noise is $1.8\%$, about four times the per-shot value $0.43\%$ for $r=d$ rounds. Each additional round lowers the threshold, because failing chains can run through more of space-time, and with ten rounds it is $0.61\%$. A preparation that is verified or consumed after a few rounds therefore tolerates two to four times more noise than a stored state.

The two-level structure of the code carries over to circuit-level decoding. The frame-deterministic checks form the upper level, which decides the logical outcome, and the gauge checks form the lower level. Hierarchical correlated matching couples the two levels with two passes of matching and raises the threshold by $22\%$ under depolarizing noise and by about $50\%$ under $Z$-biased noise. Belief-matching raises the threshold under strong bias further, to $0.86\%$, twice that of matching, at a cost per shot that is three orders of magnitude higher in our implementation. Passing hard erasures from the link level to the upper level performs worse than matching alone. The tolerance of the code to biased noise therefore appears at the circuit level only when the decoder reads the gauge checks. The coset free-energy decoder goes one step further and compares the two logical classes instead of two candidate errors. Its two-class search, local descent and entropy correction make it as accurate as Tesseract at $d=5$. Under strongly biased noise it fails on $11\%$ to $28\%$ fewer shots than belief-matching, with a gain that grows with the distance. Under depolarizing noise the gain disappears at $d=7$, where the ordered-statistics search no longer finds the best candidates. On hardware with native pair measurements, the links are measured directly and the plaquettes through cat states of ancillas. This extraction has fault distance $d$ and a threshold of $0.28\%$ with hierarchical correlated matching under the EM3 model.

We built one-parameter noise models of Helios and H2 from the rates of their data sheets, with the two-qubit error rate as the parameter, and checked them against independent measurements and against the vendor's emulator, which treats the measurement crosstalk in the same way. With these rates a single round prepares $\plusL$ with a logical error rate of $8.5\times10^{-4}$ at $d=3$ and $8.4\times10^{-5}$ at $d=5$, using 35 and 95 qubits. The $d=5$ value is six times below the physical SPAM error, so this experiment would demonstrate a logical preparation and readout that outperform their physical counterparts. Dephasing during transport limits these numbers. The crosstalk of mid-circuit measurements grows with the number of ancillas, so the crossing point of consecutive distances falls as the code grows, below the operating point of Helios between $d=7$ and $d=9$ when the state is stored for $d$ rounds. Without the crosstalk the crossings stop moving, and the model has a threshold at about three times the operating point.

Several directions remain open. Exchanging the roles of the two classes of checks gives a second frame, which we expect to prepare an eigenstate of a second logical operator. This would extend the protocol to the other logical states needed for computation. A tensor-network decoder~\cite{chubb2021} adapted to circuit-level noise would show how close the coset free-energy decoder comes to the optimal decoder at larger distances. Its candidate search needs to scale with the code to keep the gain under depolarizing noise, and a compiled implementation would make it fast enough for real-time use. Finally, the measurements of a hardware run determine the detector error model directly, which can replace the data-sheet noise models used here.

\section*{Data availability}

The circuits, decoders, simulation data and analysis scripts that support the findings of this article are available in the accompanying repository~\cite{repo}. We simulate the circuits with Stim~1.16~\cite{gidney2021stim}, match with PyMatching~2.4~\cite{higgott2022pymatching,higgott2025sparse}, run belief propagation and ordered-statistics decoding with the \texttt{ldpc} package~\cite{roffe2020} and use Tesseract~\cite{beni2025} as the reference decoder. The repository also exports every circuit as OpenQASM~2.0 for the Quantinuum processors.
